\chapter{观察性研究}

再回到那个因果问题：“一个人抬头看天，会让其他行人跟着抬头看吗？”按照上一章的做法考虑随机化实验之后，你得出结论：抬头看那么多次太费时间，而且对你的颈骨不好。于是你决定做下面这项研究：在附近找一个站在街角、没有抬头的行人；再找一个朝第一个人走来、同样没有抬头的行人。观察并记录他们在接下来 10 秒内的行为。把这一过程重复几千次。现在你就可以比较：在第一人抬头之后才抬头的第二人所占的比例，与在第一人抬头之前就已抬头的第二人所占的比例。这种由研究者观察并记录相关数据的科学研究，称为\emph{观察性研究}。

如果你真做了上述观察性研究，批评者可以说：两个行人也许都抬头了，但原因并不是第一个人的抬头引起了第二个人的抬头，而是他们都听到了头顶的一声闷雷，或者开始落雨点；因此，你的研究结果无法判定一个人的抬头是否会使他人抬头。这些批评不适用于随机化实验——这也是随机化实验在因果推断理论中居于核心地位的原因之一。然而在实践中，随机化实验对于估计因果效应的重要性要有限得多。许多科学研究并不是实验。人类的许多知识来自观察性研究。想想进化论、板块构造、全球变暖，或者天体物理学。想想人类是如何知道热咖啡可能造成烫伤的。本章回顾观察性研究能够得出有效因果推断的若干条件。

\section{可识别条件}

理想的随机化实验可用于识别并量化平均因果效应，因为随机分配处理带来了可交换性。以心脏移植与死亡率的边缘随机化实验为例：如果那些接受了移植的人没有接受移植，本应期望他们的死亡风险与那些实际未接受心脏移植的人相同。因此，该随机化实验中 0.7 的关联风险比，预期等于因果风险比。

\mnote{为简单起见，本章只考虑所有参与者始终处于随访之中、并自始至终遵守所分配治疗的随机化实验。第 8 章与第 9 章讨论其他情形。}

相比之下，观察性研究可能远没有那么令人信服（例子见本章开篇）。我们不愿赋予观察性关联以因果解释，一个关键原因就是缺乏随机的处理分配。以一项心脏移植与死亡率的观察性研究为例，其中接受心脏移植的人更可能患有严重的心脏疾病。于是，如果那些接受移植的人没有接受移植，本应期望他们的死亡风险高于那些实际未接受心脏移植的人。因此，该观察性研究中 1.1 的关联风险比，是两种力量妥协的结果：移植对死亡率的真实有益效应（它把关联风险比推向 1 以下），以及接受移植者本已更高的死亡风险（它把关联风险比推向 1 以上）。观察性研究中处理与结局之间的关联，最好的解释未必是处理对结局的因果效应。

尽管承认随机化实验在因果推断上具有内在优势，有时我们仍不得不依靠观察性研究来回答因果问题。那该怎么办？一种常见策略是：把数据当作处理是在给定已测量协变量 $L$ 的条件下随机分配的来分析——尽管我们往往知道这充其量只是一种近似。于是，从观察性数据做因果推断，就围绕着这样一个希望：该观察性研究可以被视为条件随机化实验。非正式地说，若下列条件成立，观察性研究就可以被概念化为条件随机化实验：

\begin{enumerate}
\item 观测到的结局等于观测到的治疗取值之下的反事实结局；
\item 反事实结局独立于观测到的处理，且仅以已知的已测量协变量 $L$ 为条件；
\item 在给定 $L$ 的条件下，接受处理每一个取值的概率都大于零，即正值性。
\end{enumerate}

本章将在观察性研究的语境中逐一展开这三个条件。条件 1 在第 1 章中称为一致性，条件 2 在前几章中称为可交换性，条件 3 在技术注记 2.3 中称为正值性。

我们将看到，这些条件往往是英雄式的，这也解释了为何观察性研究得出的因果推断常被怀疑。不过，如果观察性研究与条件随机化实验之间的类比恰好成立，我们就可以用上一章描述的方法——逆概率加权或标准化——从观察性研究中识别因果效应。因此，我们把这些条件称为\emph{可识别条件}或\emph{可识别假设}。例如，在上一章中，我们用表 2.2 的数据算得因果风险比等于 1，而那批数据来自一项条件随机化实验。如果同一批数据（现在列于表 3.1）来自一项观察性研究，且上述三个可识别条件成立，我们同样会算得因果风险比等于 1。

\mnote{表 3.1}

重要的是，在理想的随机化实验中，可识别条件由设计本身保证成立。也就是说，对条件随机化实验，我们只需要表 3.1 的数据就能算出因果风险比 1。相反，要从观察性研究识别因果风险比，我们必须假定可识别条件成立，而这一假定当然可能不真。从观察性数据做因果推断需要两样东西：数据，以及可识别条件。关于可识别性更精确的定义，见精要 3.1。

\mnote{Rubin (1974, 1978) 把 Neyman 关于随机化实验的理论推广到观察性研究。Rosenbaum and Rubin (1983) 把可交换性与正值性的组合称为弱可忽略性，把完全可交换性（见技术注记 2.1）与正值性的组合称为强可忽略性。}

\begin{table}[htbp]
\centering
\caption{心脏移植研究的数据：$L$ 为基线危重状态（1：危重），$A$ 为心脏移植（1：接受），$Y$ 为 5 年内死亡（1：死亡）}
\label{tab:3.1}
\small
\begin{tabular}{lccc@{\hspace{2.5em}}lccc}
\toprule
& $L$ & $A$ & $Y$ & & $L$ & $A$ & $Y$\\
\midrule
Rheia      & 0 & 0 & 0 & Leto       & 1 & 0 & 0\\
Kronos     & 0 & 0 & 1 & Ares       & 1 & 1 & 1\\
Demeter    & 0 & 0 & 0 & Athena     & 1 & 1 & 1\\
Hades      & 0 & 0 & 0 & Hephaestus & 1 & 1 & 1\\
Hestia     & 0 & 1 & 0 & Aphrodite  & 1 & 1 & 1\\
Poseidon   & 0 & 1 & 0 & Polyphemus & 1 & 1 & 1\\
Hera       & 0 & 1 & 0 & Persephone & 1 & 1 & 1\\
Zeus       & 0 & 1 & 1 & Hermes     & 1 & 1 & 0\\
Artemis    & 1 & 0 & 1 & Hebe       & 1 & 1 & 0\\
Apollo     & 1 & 0 & 1 & Dionysus   & 1 & 1 & 0\\
\bottomrule
\end{tabular}
\end{table}

当任一可识别条件不成立时，观察性研究与条件随机化实验之间的类比就失效了。在这种情况下，从观察性数据做因果推断还有其他可能途径，它们需要另一组可识别条件。其中一种途径是寄望于某个处理预测因子——称为\emph{工具变量}——表现得就像是在给定已测量协变量的条件下被随机分配的一样。我们在第 16 章讨论工具变量方法。

\renewcommand{\FPname}{精要（Fine Point）}
\begin{finepoint}[因果效应的可识别性]
若某一组假设意味着观测数据的分布与效应度量的某一个取值相容，我们就说平均因果效应在这组假设下是（非参数）可识别的。反之，若观测数据的分布与效应度量的若干取值都相容，我们就说平均因果效应在这些假设下不可识别。例如，如果表 3.1 的研究来自一项条件随机化实验，其中接受处理的概率依赖于 $L$ 的取值（从而条件可交换性 $Y^{a}\indep A\given L$ 由设计成立），那么上一章已表明因果效应可识别：因果风险比等于 1，且无需任何进一步假设。然而，如果表 3.1 的数据来自观察性研究，那么只有当我们在数据之外补充条件可交换性 $Y^{a}\indep A\given L$ 这一假设时，因果风险比才等于 1。要在观察性研究中识别因果效应，我们需要一个外在于数据的假设，即识别假设。事实上，如果我们决定不为数据补充识别假设，那么表 3.1 的数据与下列因果风险比都相容：

\begin{itemize}
\item 小于 1，若 $L$ 之外的风险因素在接受治疗者中更常见；
\item 大于 1，若 $L$ 之外的风险因素在未接受治疗者中更常见；
\item 等于 1，若除 $L$ 之外的一切风险因素在接受治疗者与未接受治疗者之间分布相同，或等价地，若 $Y^{a}\indep A\given L$。
\end{itemize}

本章讨论非参数识别平均因果效应的三个可识别条件。在第 16 章，我们描述另一组足以非参数识别平均因果效应的可识别条件。
\end{finepoint}
\renewcommand{\FPname}{精要}

毫不奇怪，基于与条件随机化实验之类比的观察性方法，在那种类比常常合理的学科（如流行病学）中历来受到偏爱；而工具变量方法则在观察性研究常常无法在给定已测量协变量下被概念化为条件随机化实验的学科（如经济学）中历来受到偏爱。在第 16 章之前，我们将集中讨论依赖观察性研究能否模拟条件随机化实验的因果推断途径。下面更详细地描述这三个可识别条件。

\section{可交换性}

关于可交换性 $Y^{a}\indep A$，我们已经谈了很多。在边缘（即无条件）随机化实验中，受治疗者与未受治疗者可交换，因为受治疗者若一直接受未治疗，本会经历与未受治疗者相同的平均结局；反之亦然。之所以如此，是因为随机化保证了结局的独立预测因子在受治疗组与未治疗组之间分布相同。

\mnote{结局的独立预测因子，是在各治疗水平之内与结局 $Y$ 相关的协变量。对二值结局，结局的独立预测因子常被称为结局的风险因素。}

以表 3.1 概括的研究为例。上一章我们说过，在这项研究中可交换性显然不成立，因为基线时受治疗者中有 69\%、未受治疗者中有 43\% 处于危重状态 $L=1$。这种独立结局预测因子分布上的不平衡，在边缘随机化实验中不应出现（实际上，这种不平衡可能因偶然出现，但让我们继续假定我们的研究足够大，不会出现偶然的发现）。

另一方面，在条件随机化实验中——即接受处理的概率依赖于 $L$ 的实验——独立结局预测因子 $L$ 在受治疗者与未受治疗者之间分布不平衡，正是设计所预期的。表 3.1 的研究就是这样的条件随机化实验：受治疗者与未受治疗者不可交换——因为受治疗者在研究开始时平均预后更差——但在变量 $L$ 的各层之内，受治疗者与未受治疗者是条件可交换的。在子集 $L=1$（危重状态）中，受治疗者与未受治疗者可交换，因为受治疗者若一直接受未治疗，本会经历与未受治疗者相同的平均结局；反之亦然。对子集 $L=0$ 同理。等价的说法是：在条件随机化实验中条件可交换性 $Y^{a}\indep A\given L$ 成立，因为在 $L$ 的各层之内，其余一切结局预测因子在受治疗组与未治疗组之间分布相同。

回到观察性研究。当处理不是由研究者随机分配时，接受治疗的原因很可能与某些结局预测因子相关。也就是说，如同条件随机化实验，观察性研究中结局预测因子的分布一般会在受治疗组与未治疗组之间不同。例如，表 3.1 的数据可能来自这样一项观察性研究：医生倾向于把稀缺的心脏移植给最需要的人，即处于危重状态 $L=1$ 的个体。事实上，如果受治疗者与未受治疗者之间分布不平等的唯一结局预测因子就是 $L$，那么表 3.1 的研究既可以称为 (i) 一项观察性研究，其中接受治疗 $A=1$ 的概率在 $L=1$ 者中为 0.75、在 $L=0$ 者中为 0.50；也可以称为 (ii) 一项（非盲）条件随机化实验，其中研究者对 $L=1$ 者以 0.75 的概率、对 $L=0$ 者以 0.50 的概率随机分配治疗 $A=1$。这两种刻画在逻辑上等价。在任一刻画之下，条件可交换性 $Y^{a}\indep A\given L$ 都成立，且可用标准化或逆概率加权识别因果效应。

\mnote{第 7 章我们把这类结局预测因子称为混杂因子。}

当然，对这项观察性研究而言，关键问题是：$L$ 是否是受治疗者与未受治疗者之间分布不平等的唯一结局预测因子？遗憾的是，这个问题注定无法回答，因此我们的研究者必须愿意在条件可交换性 $Y^{a}\indep A\given L$ 成立的假设下工作。另外要注意，并非所有在治疗组之间分布不平等的变量都需要纳入 $L$。例如，心脏移植会被分配给排斥移植概率低的个体，即带有某些人类白细胞抗原（HLA）基因的心脏会被分配给恰好拥有相容基因的个体。由于 HLA 基因不是死亡率的预测因子（在给定 $L$ 与 $A$ 的条件下），治疗分配在 $L$ 的各层之内实质上是随机的，因而 HLA 无须在分析中考虑。

在没有随机化的情况下，无法保证条件可交换性成立。例如，设想研究者并不知道，医生更愿意把心脏移植给不吸烟者。考虑两个 $L=1$ 的个体：一人吸烟（$U=1$），另一人不吸烟（$U=0$）；$U=1$ 的那位接受治疗 $A=1$ 的概率更低。当吸烟这一重要的结局预测因子的分布在 $L=1$ 层内受治疗者（吸烟者 $U=1$ 的比例较低）与未受治疗者（吸烟者比例较高）之间不同时，给定 $L$ 的条件可交换性不成立。重要的是，即使收集了吸烟的数据，也无法排除其他未知的、分布不平衡的结局预测因子仍未被测量。

\mnote{我们用 $U$ 表示未测量的变量。由于未测量的变量不能用于标准化或逆概率加权，当已测量变量 $L$ 不足以实现条件可交换性时，因果效应无法识别。}

因此，在观察性研究中可交换性 $Y^{a}\indep A\given L$ 可能不成立。具体而言，若存在未测量的结局独立预测因子 $U$，使得在 $L$ 的各层之内接受治疗 $A$ 的概率依赖于 $U$，则条件可交换性 $Y^{a}\indep A\given L$ 不成立。更糟的是，即便条件可交换性 $Y^{a}\indep A\given L$ 成立，研究者也无法经验地核实它确实成立。如果他们根本没有收集吸烟的数据，又怎能检查吸烟的分布在受治疗者与未受治疗者中相同呢？那些同样可能在受治疗者与未受治疗者之间分布不同的其他未测量结局预测因子 $U$ 又该怎么办？在条件可交换性下分析观察性研究时，我们只能希望自己的专家知识能正确引导我们收集足够的数据，使该假设至少近似成立。

研究者可以运用专家知识来提高条件可交换性假设的可信度。他们可以测量许多相关变量 $L$（例如既是处理的决定因素、又是独立结局预测因子的变量），而不像表 3.1 那样只测一个变量，然后假定条件可交换性在由所有这些变量 $L$ 的组合所定义的分层之内近似成立。不幸的是，无论 $L$ 中纳入多少变量，都无法检验该假设是否正确，这使从观察性数据做因果推断成为一项有风险的任务。因果推断的有效性要求研究者的专家知识是正确的。这些知识以“给定已测量协变量的可交换性”这一假设的形式编码，为数据提供补充，以期识别所关心的因果效应。

\mnote{要验证条件可交换性，需比较反事实 $Y^{a}$ 在 $A=a$ 与 $A\neq a$ 两组（给定 $L$）的分布是否相同。但这在逻辑上不可能：对未接受治疗 $a$ 的个体（$A\neq a$），$Y^{a}$ 的取值未知，故无从作此比较。}

\renewcommand{\FPname}{精要（Fine Point）}
\begin{finepoint}[交叉随机化实验]
在精要 2.1 中，我们描述了交叉实验：一个个体的观测跨越两个或更多时段——比如 $t=0$ 与 $t=1$——且该个体在每个时段接受不同的处理取值。我们曾表明，当下列三个很强的条件成立时，交叉实验中的个体因果效应可识别：i) 处理无延续效应：$Y_{it=1}^{a_0,a_1}=Y_{it=1}^{a_1}$；ii) 个体因果效应不随时间变化：$Y_{it}^{a_t=1}-Y_{it}^{a_t=0}=\alpha_i$，$t=0,1$；iii) 无处理之下的反事实结局不随时间变化：$Y_{it}^{a_t=0}=\beta_i$，$t=0,1$。这不需要任何随机化。现在我们转向交叉随机化实验，其中个体接受处理取值的顺序是被随机分配的。

当由于可能存在时间效应而不再假定 iii) 成立时，随机的处理分配就变得重要了。为简单起见，假定每个个体以 0.5 的概率被随机分配到 $(A_{i1}=1,A_{i0}=0)$ 或 $(A_{i1}=0,A_{i0}=1)$。令 $Y_{i0}^{a_1=0}-Y_{i0}^{a_0=0}=r_i$。那么在 i) 与 ii) 以及一致性之下，若 $A_{i0}=0$ 且 $A_{i1}=1$，则 $Y_{i1}-Y_{i0}=\alpha_i+r_i$；若 $A_{i1}=0$ 且 $A_{i0}=1$，则 $Y_{i0}-Y_{i1}=\alpha_i-r_i$。由于 $r_i$ 未知，我们不再能识别个体因果效应；但由于 $A_{i1}$ 与 $A_{i0}$ 是随机的、因而独立于 $r_i$，$(Y_{i1}-Y_{i0})A_{i1}+(Y_{i0}-Y_{i1})A_{i0}$ 的均值估计了平均因果效应，即 $\E[\alpha_i]$。若我们只假定 i)，则该均值估计的是时刻 0 与时刻 1 平均因果效应的平均，即 $(\E[\alpha_{i1}]+\E[\alpha_{i0}])/2$，其中 $\alpha_{it}=Y_{it}^{a_t=1}-Y_{it}^{a_t=0}$。

总之，若无延续效应这一假设 1) 成立，交叉实验就可用来估计平均因果效应。然而，对于本书所研究的处理与结局类型，无延续效应这一假设并不合理。
\end{finepoint}
\renewcommand{\FPname}{精要}

\section{正值性}

有些研究者计划做一项实验，以计算心脏移植 $A$ 对 5 年死亡率 $Y$ 的平均效应。不言而喻，研究者会把一些个体分配到治疗水平 $A=1$，另一些分配到 $A=0$。设想另一种做法：研究者把所有个体都分配到 $A=1$ 或都分配到 $A=0$。那太荒唐了。若所有人都接受同一治疗水平，计算平均因果效应将无从谈起。因此，我们必须这样分配处理：使几乎可以肯定，每个治疗组都会分到一些个体。换言之，我们必须确保被分配到每个治疗水平的概率都大于零——是一个正概率。这就是\emph{正值性}条件。

\mnote{正值性条件有时被称为实验性治疗假设。}

在描述实验时我们没有强调正值性，因为在这类研究中正值性被视为理所当然。在边缘随机化实验中，$\Pr[A=1]$ 与 $\Pr[A=0]$ 按设计都为正。在边缘与条件随机化实验中，对变量 $L$ 的一切水平 $l$，条件概率 $\Pr[A=1\given L=l]$ 与 $\Pr[A=0\given L=l]$ 也按设计为正。例如，如果表 3.1 的数据来自条件随机化实验，那么被分配到心脏移植的条件概率在危重者中为 $\Pr[A=1\given L=1]=0.75$，在其他人中为 $\Pr[A=1\given L=0]=0.50$。给定 $L$ 时正值性成立，因为这两个概率都不为 0（也不为 1，否则意味着不接受心脏移植 $A=0$ 的概率为 0）。因此我们说，若对所涉因果对比中的每个 $a$ 都有 $\Pr[A=a\given L=l]>0$，则正值性成立。实际上，这个定义并不完整：如果我们的研究总体仅限于 $L=1$ 组，那就不需要在 $L=0$ 组中要求正值性。正值性只对所关心总体中实际存在的那些 $l$ 取值有要求。

\mnote{\parbox[t]{\marginparwidth}{\raggedleft 正值性：$\Pr[A=a\given L=l]>0$，对所有满足 $\Pr[L=l]\neq 0$ 的 $l$。}}

此外，正值性只对可交换性所需的那些变量 $L$ 有要求。例如，在表 3.1 的条件随机化实验中，我们不会去问蓝眼睛的人接受治疗的概率是否大于 0，因为“蓝眼睛”这一变量对实现受治疗者与未受治疗者之间的可交换性并不必要。（在给定 $L$ 与 $A$ 的条件下，“蓝眼睛”不是结局 $Y$ 的独立预测因子，甚至未被用于分配治疗。）也就是说，标准化风险与逆概率加权风险在仅对 $L$ 调整之后就等于反事实风险；对那些像“蓝眼睛”一样无须调整的变量，并不需要正值性。

在观察性研究中，正值性与可交换性都没有保证。例如，如果医生总是把心脏移植给危重状态的个体 $L=1$，即 $\Pr[A=0\given L=1]=0$，那么正值性不成立，如图 3.1 所示。可交换性与正值性这两个条件的一个差别在于：正值性有时可以被经验地核实（见第 12 章）。例如，如果表 3.1 对应的是观察性研究的数据，我们会断定正值性对 $L$ 成立，因为在 $L$ 的每个水平（即 $L=0$ 与 $L=1$）上，治疗的各个水平（即 $A=0$ 与 $A=1$）都有人。我们在上一章对标准化与逆概率加权的讨论明确提到了可交换性条件，却只是隐含地假定了正值性条件（技术注记 2.3 中则是明确的）。我们此前对标准化风险与逆概率加权风险的定义，实际上只有在正值性成立时才有意义。为了直观理解为什么正值性条件失效时标准化风险与逆概率加权风险没有良好定义，请看图 3.1。如果不存在 $L=1$ 的未受治疗者（$A=0$），数据中就没有任何信息可以用来模拟“所有受治疗者都未受治疗”会怎样，因为不存在可与 $L=1$ 的受治疗者相交换的 $L=1$ 未受治疗者。细节见技术注记 3.1。

\mnote{图 3.1}

\renewcommand{\TPname}{技术注记（Technical Point）}
\begin{techpoint}[标准化与逆概率加权下的正值性]
我们已把治疗水平 $a$ 的标准化均值定义为 $\sum_l \E[Y\given A=a,L=l]\Pr[L=l]$。然而，只有当条件量 $\E[Y\given A=a,L=l]$ 有良好定义时，这个表达式才能计算；而当条件概率 $\Pr[A=a\given L=l]$ 对总体中出现的一切 $l$ 取值都大于零时，情况正是如此。也就是说，当正值性成立时。（注意，“对一切满足 $\Pr[L=l]\neq 0$ 的 $l$ 都有 $\Pr[A=a\given L=l]>0$”这一陈述，实际上等价于“以概率 1 有 $f[a\given L]>0$”。）因此，标准化均值定义为

\[
\begin{gathered}
\sum_l \E[Y\given A=a,L=l]\Pr[L=l]\\
\text{若对一切满足 } \Pr[L=l]\neq 0 \text{ 的 } l \text{ 都有 } \Pr[A=a\given L=l]>0,
\end{gathered}
\]

否则无定义。只有当对总体中协变量 $L$ 的每个取值，都有一些个体接受了治疗水平 $a$ 时，标准化均值才能计算。

当正值性不成立时，逆概率加权均值 $\E\!\left[\frac{I(A=a)Y}{f(A\given L)}\right]$ 不再等于 $\E\!\left[\frac{I(A=a)Y}{f(a\given L)}\right]$。具体而言，$\E\!\left[\frac{I(A=a)Y}{f(a\given L)}\right]$ 无定义，因为计算期望时会出现无定义的比值 $0/0$。另一方面，逆概率加权均值 $\E\!\left[\frac{I(A=a)Y}{f(A\given L)}\right]$ 总有良好定义，因为其分母 $f[A\given L]$ 永远不可能为零。然而，当正值性不成立时，即便在可交换性之下，它现在也是反事实均值的有偏估计。特别地，$\E\!\left[\frac{I(A=a)Y}{f(A\given L)}\right]$ 等于

\[
\Pr[L\in Q(a)]\sum_l \E[Y\given A=a,L=l,L\in Q(a)]\Pr[L=l\given L\in Q(a)],
\]

其中 $Q(a)=\{l;\Pr(A=a\given L=l)>0\}$ 是那些使 $A=a$ 可以以正概率被观测到的 $l$ 取值之集合。因此，在可交换性之下，$\E\!\left[\frac{I(A=a)Y}{f(A\given L)}\right]$ 等于 $\E[Y^{a}\given L\in Q(a)]\Pr[L\in Q(a)]$。

由 $Q(a)$ 的定义，当 $A$ 为二值且正值性不成立时，$Q(0)$ 不可能等于 $Q(1)$。此时对比 $\E\!\left[\frac{I(A=1)Y}{f(A\given L)}\right]-\E\!\left[\frac{I(A=0)Y}{f(A\given L)}\right]$ 没有因果解释，即便在可交换性之下也是如此，因为它是两个不同组之间的对比。当正值性成立时，$Q(1)=Q(0)$，且若可交换性成立，该对比就是平均因果效应。
\end{techpoint}
\renewcommand{\TPname}{技术注记}

\begin{figure}[htbp]
\centering
\begin{tikzpicture}[causal,scale=0.9,every node/.style={transform shape},
    lab/.style={font=\scriptsize,inner sep=1.2pt}]
  % 分支连线与标注（先画线，节点随后覆盖端点）
  \draw[ed] (0,0) -- (6.4,4.2)
     node[pos=0.45,sloped,above=1pt,fill=white,inner sep=1pt,font=\footnotesize] {$A\!=\!0\ (.5)$}
     node[pos=0.60,sloped,below=1pt,fill=white,inner sep=0.5pt,font=\footnotesize] {4};
  \draw[ed] (0,0) -- (6.4,1.6)
     node[pos=0.45,sloped,above=1pt,fill=white,inner sep=1pt,font=\footnotesize] {$A\!=\!1\ (.5)$}
     node[pos=0.60,sloped,below=1pt,fill=white,inner sep=0.5pt,font=\footnotesize] {4};
  \draw[ed] (0,0) -- (6.4,-1.6)
     node[pos=0.45,sloped,above=1pt,fill=white,inner sep=1pt,font=\footnotesize] {$A\!=\!0\ (0)$}
     node[pos=0.60,sloped,below=1pt,fill=white,inner sep=0.5pt,font=\footnotesize] {0};
  \draw[ed] (0,0) -- (6.4,-4.2)
     node[pos=0.45,sloped,above=1pt,fill=white,inner sep=1pt,font=\footnotesize] {$A\!=\!1\ (1)$}
     node[pos=0.60,sloped,below=1pt,fill=white,inner sep=0.5pt,font=\footnotesize] {12};
  % 根节点
  \draw[fill=white,thick] (0,0) circle (1.0);
  \draw[thick] (-0.90,-0.24) -- (0.90,0.24);
  \node[lab] at (0,0.46) {$L\!=\!0\ (.4)$};
  \node[lab] at (0,-0.46) {$L\!=\!1\ (.6)$};
  % 叶节点 1
  \draw[fill=white,thick] (6.4,4.2) circle (1.0);
  \draw[thick] (5.50,3.96) -- (7.30,4.44);
  \node[lab] at (6.4,4.65) {$Y\!=\!0\ (.75)$};
  \node[lab] at (6.4,3.75) {$Y\!=\!1\ (.25)$};
  \node[font=\footnotesize] at (7.80,4.60) {3};
  \node[font=\footnotesize] at (7.80,3.80) {1};
  % 叶节点 2
  \draw[fill=white,thick] (6.4,1.6) circle (1.0);
  \draw[thick] (5.50,1.36) -- (7.30,1.84);
  \node[lab] at (6.4,2.05) {$Y\!=\!0\ (.75)$};
  \node[lab] at (6.4,1.15) {$Y\!=\!1\ (.25)$};
  \node[font=\footnotesize] at (7.80,2.00) {3};
  \node[font=\footnotesize] at (7.80,1.20) {1};
  % 叶节点 3
  \draw[fill=white,thick] (6.4,-1.6) circle (1.0);
  \draw[thick] (5.50,-1.84) -- (7.30,-1.36);
  \node[lab] at (6.4,-1.15) {$Y\!=\!0$};
  \node[lab] at (6.4,-2.05) {$Y\!=\!1$};
  \node[font=\footnotesize] at (7.80,-1.20) {0};
  \node[font=\footnotesize] at (7.80,-2.00) {0};
  % 叶节点 4
  \draw[fill=white,thick] (6.4,-4.2) circle (1.0);
  \draw[thick] (5.50,-4.44) -- (7.30,-3.96);
  \node[lab] at (6.4,-3.75) {$Y\!=\!0\ (1/3)$};
  \node[lab] at (6.4,-4.65) {$Y\!=\!1\ (2/3)$};
  \node[font=\footnotesize] at (7.80,-3.80) {4};
  \node[font=\footnotesize] at (7.80,-4.60) {8};
\end{tikzpicture}
\caption{正值性失效：在 $L=1$ 层内无人未受治疗（$\Pr[A=0\given L=1]=0$），因而没有可与该层受治疗者相交换的未受治疗者，以模拟“所有受治疗者都未受治疗”会怎样。括号内为相应分支上的条件概率，分支旁与右侧的数字为个体数}
\label{fig:3.1}
\end{figure}

\section{一致性：首先，定义反事实结局}

反事实的一致性意味着：每个受治疗个体的观测结局 $Y$ 等于她若接受治疗本会有的结局 $Y^{a=1}$；每个未受治疗个体的观测结局 $Y$ 等于她若一直接受未治疗本会有的结局 $Y^{a=0}$。也就是说，一致性是这样一种假设：$Y=Y^{A}$，其中 $Y^{A}$ 是把反事实 $Y^{a}$ 中的 $a$ 取为个体实际治疗 $A$ 的结果。

一致性是因果推断的基本条件，因为它把反事实 $Y^{a}$ 与观测数据 $Y$ 联系起来。在本书中，我们把反事实 $Y^{a}$ 当作初始概念。也就是说，观测结局 $Y$ 由这些初始概念通过公式 $Y=Y^{A}$ 导出（即是其函数）。对二值处理 $A$ 很容易看出，$Y^{A}$ 只依赖于初始概念，因为此时 $Y^{A}=AY^{a=1}+(1-A)Y^{a=0}$。

\mnote{Robins and Greenland (2000) 主张，定义良好的反事实（或数学上等价的构造）是有意义因果推断的必要条件。}

在某些情形下，一致性看起来显然成立。例如，如果你吃了一片阿司匹林 $A=1$ 而后死亡（$Y=1$），那么你若吃了阿司匹林的反事实结局 $Y^{a=1}$ 必定等于 1。但一致性并非总能被视为理所当然，理由见下文。

一致性条件有两个主要组成部分：(1) 通过指定上标 $a$ 来精确地定义反事实结局 $Y^{a}$；(2) 把反事实结局与观测结局联系起来。本节处理一致性的第一个组成部分。

本书描述的因果推断方法，以定义良好的反事实结局 $Y^{a}$ 之存在为前提。显然，一个个体的因果效应 $Y^{a=1}-Y^{a=0}$ 只有在 $Y^{a}$ 对 $a=1$ 与 $a=0$ 都有良好定义时才有良好定义。由此可知，对给定总体，平均因果效应 $\E[Y^{a=1}-Y^{a=0}]=\Pr[Y^{a=1}=1]-\Pr[Y^{a=0}=1]$ 只有在总体中每个个体的个体因果效应都有良好定义时才有良好定义。于是关键问题就是：“我们怎么知道反事实有良好定义？”一个自然且可取的充分条件是：若 $a$ 对应一个定义良好的干预，则 $Y^{a}$ 有良好定义，即“在实施干预 $a$ 之下的结局”。为说明定义良好的干预这一概念，考虑两项理想的随机化实验（其中所有个体都接受了分配给他们的治疗），它们在同一总体中的个体身上进行，参与者被随机分配到心脏移植 $a=1$ 或药物治疗 $a=0$。总体中所有个体都有资格接受 $a=1$ 或 $a=0$。

在第一个随机化实验中，研究者写了一份方案，详细描述了两个所关心的干预。研究者规定，被分配到心脏移植的个体要接受某些术前程序、麻醉、手术技术、术后护理与免疫抑制治疗，以期确保每个被分配到心脏移植的个体接受相同的治疗 $a=1$；对 $a=0$ 也作类似规定。如果方案没有规定这些细节，那么每位医生可能实施了不同版本的“心脏移植”，也许采用各自偏爱的手术技术或免疫抑制治疗。在这项研究中，“心脏移植”一词对应一个定义良好的干预 $a=1$，它被定义为对所有个体依次实施各组成部分：$a_0=1$（分配到心脏移植）、$a_1$（预先规定的术前程序）、$a_2$（麻醉）与 $a_3$（手术技术）。

在第二个随机化实验中，研究者有意不对干预作精确规定，以便试验中实施的干预能反映真实世界中的情形。在这项研究中，“心脏移植”一词对应一个定义良好的干预 $a=1$，它被定义为“分配到心脏移植”（即对所有个体 $a_0=1$），随后观察干预之后自然发生的一切。也就是说，对每个被分配到心脏移植的个体，$a_1$、$a_2$、$a_3$ 的取值并不事先规定，而是医疗系统中自然出现的取值：对每个个体 $a_1=A_1$、$a_2=A_2$、$a_3=A_3$。

\mnote{这类试验常被称为实用性试验。}

这两个例子说明了实践中常见的一种情形：同一个治疗名称被用于不同的含义。在第一个实验中，“心脏移植”这一标签 $a=1$ 指的是依次实施各组成部分干预 $(a_0,a_1,a_2,a_3)$。在第二个实验中，它指的是实施时点干预 $a_0$，随后研究者让世界自行其是。\mnote{形式上，第一个实验的 $Y^{a=1}$ 是联合反事实 $Y^{a_0=1,a_1,a_2,a_3}$；第二个实验的则把 $a_1,a_2,a_3$ 换为自然取值 $A_1,A_2,A_3$，即 $Y^{a_0=1,A_1,A_2,A_3}$。对第二个实验中分配到心脏移植 $A=1$ 者，它等于观测结局 $Y$。见技术注记 3.2。}因此，$(A_1,A_2,A_3)$ 的取值将取决于总体的特征以及实验所处的环境。

但是，尽管每项实验实施的都是 $a=1$ 的不同版本，相应的干预 $a=1$ 在每项实验的方案中都是定义良好的。这意味着，反事实结局 $Y^{a=1}$ 在每项实验中都有良好定义：若该实验方案中关于干预 $a=1$ 的指示被遵循，个体本会有的结局（$Y^{a=0}$ 同理）。然而，两项实验的反事实结局 $Y^{a}$ 很可能不同。若果真如此，那么因果效应 $\Pr[Y^{a=1}=1]-\Pr[Y^{a=0}=1]$ 虽然在每项实验中都有良好定义，在两项实验之间也会不同。这就引出了一个问题：研究成果的使用者更偏爱这两个因果效应中的哪一个？

第一个实验的因果效应在不同总体之间的变化幅度，可能小于第二个实验的因果效应（因为前者精确规定了干预的各个组成部分，而后者的干预后变量之自然取值可能因总体而异）。因此，在某一特定总体中进行的第一个实验所得的因果效应，也许比第二个实验的因果推断更容易迁移到其他总体。注意，这里讨论的是因果效应的\emph{可迁移性}，而不是各项研究中因果效应是否有良好定义。

\mnote{第 4 章讨论可能影响因果效应可迁移性的若干因素。}

当干预对应于可以作为实验方案一部分加以描述的行动时，它就是定义良好的。然而，定义良好的干预不必被完美地规定。事实上，完美地规定干预通常是不可能的。再考虑第一个实验。其方案规定了组成部分 $a_0$、$a_1$、$a_2$、$a_3$，却没有规定实施心脏移植的外科医生的训练。因此，经验丰富与缺乏经验的外科医生都可能参与了研究。由于已知移植经验不足会影响移植后死亡率，若所有人都按方案接受治疗，风险 $\Pr[Y^{a=1}=1]$ 将取决于参与研究的外科医生经验的未知分布。即便实验规定了外科医生的训练，我们也总能找到方案中其他未规定、或可作多种解释的东西。

干预无法被完美规定，是平均因果效应 $\Pr[Y^{a=1}=1]-\Pr[Y^{a=0}=1]$ 的取值可能因总体而异的原因之一。例如，在我们的心脏移植实验中，一个外科经验分布不同的新社区里的平均因果效应，通常会有别于试验总体中的效应，即使新总体严格遵循与试验相同的方案。我们对干预规定得越精确，我们的因果问题就越精确，一般也越容易把因果推断从一个总体迁移到另一个总体。当然，对结局没有影响的干预组成部分不会影响可迁移性。例如，在心脏移植实验中，我们不必担心外科医生手术服的颜色（绿色还是蓝色），因为科学家们一致认为改变手术服颜色不会导致不同的结局。

\mnote{“没有操纵就没有因果”（Holland 1986）这句话常被用来表达这样一种观点：有意义的因果推断需要定义良好的干预。但请记住，定义良好的干预未必是人所能实施、或在历史上某一时刻切实可行的干预。例如，在有基因修饰技术之前，基因变异对疾病的效应就被认为是定义良好的。}

上述一切考虑对随机化实验与观察性研究同样适用。无论数据如何产生，定义良好的干预 $a$ 都意味着定义良好的反事实 $Y^{a}$。考虑一项观察性研究，研究在“所有个体若被提供心脏移植都会接受”的总体中“被提供心脏移植”的效应。我们可以把干预 $a=1$ 规定为一串精确定义的组成部分（如上面的第一个实验），也可以规定为反映真实世界情形的最小化干预（如第二个实验）。

随机化实验与观察性研究的一个常见差别，在于人们对定义因果效应的干预之共识程度。在理想实验中，干预在方案中被规定并实际实施，因此研究者对干预是什么有共同理解，从而使反事实 $Y^{a}$ 有良好定义。在观察性研究中，由于干预未必在真实世界中被预先规定，研究者对何为所关心的干预可能有不同看法。例如，谈论“心脏移植”$a=1$ 的效应却未明确定义所欲干预的研究者，可能指的是各种不同的因果效应，例如一串精确定义的组成部分的效应，或反映真实世界情形的最小化干预的效应。结果，反事实 $Y^{a}$ 仍是定义不良的，“心脏移植”的因果效应也就成了一个过于模糊的概念。不过，这种模糊性有时可以通过让我们的专家研究者每次只聚焦于一个特定干预而克服。此时，每个干预都可能被视为定义良好，各自有其不同但明确的反事实 $Y^{a}$。

研究者一致认为，当就一切实际目的而言，反事实 $Y^{a}$ 不再残留有意义的模糊性时，某个特定干预 $a$ 就是\emph{足够定义良好的}。这就引出了一个问题：“我们怎么知道一个干预对我们的目的而言足够定义良好？”或者等价地问：“我们怎么知道不再残留有意义的模糊性？”答案是：“我们不知道。”宣布一个干预足够定义良好，是一群专家基于历史某一时刻可获得的实质知识所达成的共识。然而，即便专家们现在同意某个干预足够定义良好，将来有了新知识，他们也可能改变看法。因此，“足够定义良好的干预”一词依赖于可获得的知识。精要 3.3 把这一讨论与先前的提法联系起来。

\renewcommand{\FPname}{精要（Fine Point）}
\begin{finepoint}[可能世界]
科学哲学家提出了基于“可能世界”概念的反事实理论（Stalnaker 1968，Lewis 1973）。反事实 $Y^{a}$ 被定义为：在个体接受了该处理的那个可能世界中——即最接近现实世界的那个世界——$Y$ 的取值。特别地，这些哲学家假定，当 $A=a$ 时 $Y^{a}=Y$，因为最接近现实世界的可能世界就是现实世界本身。因此，在他们对反事实的定义之下，一致性总是成立。

当 $A\neq a$ 时，“最接近的可能世界”因而反事实 $Y^{a}$ 总有些定义不良、模糊不清。尽管如此，Lewis 指出他的反事实定义常常有用。Robins and Greenland (2000) 同意这一点，但还主张，应当用定义良好的干预这一概念取代最接近的可能世界这一概念，因为在观察性研究中，反事实之所以模糊、定义不良，恰恰是因为人们未能把所考虑的假想干预与因果对比精确化。
\end{finepoint}
\renewcommand{\FPname}{精要}

在观察性研究中，一个常见问题出现在研究者想要量化那些并不对应于某人实际实施过的干预的变量之因果效应时。这类非干预变量的一些例子是生物学状态（如血压、LDL 胆固醇、体重）或社会因素（如社会经济地位），它们只能通过干预其成因来改变。例如，考虑“变得肥胖对心肌梗死的影响”。这句话没有有意义的解释，因为反事实结局定义不良。人们也许以为，如果我们规定了干预的起止（如 40 岁到 50 岁）以及改变该状态或因素的程序（如药物、手术、饮食、运动），这些效应就会有良好定义。但是，如果我们规定了所有这些细节，我们所描述的就将是我们所规定的那些干预的效应，而不是（比如说）肥胖的效应。

\mnote{Hernán and Taubman (2008) 讨论了两位世界领袖——一位专制君主与一位糊涂总统——在考虑“肥胖的效应”时的困境。}

\mnote{效应是否定义不良取决于结局。考虑肥胖对求职歧视的效应——以雇主审阅求职者简历与照片后通知面试的比例来衡量。由于处理是“雇主所感知的肥胖”，导致肥胖的机制可能无关紧要。}

然而，现在考虑“血压 $A$ 对卒中 $Y$ 的效应”。由于血压的改变只能由影响血压的特定干预（如不同类型的降压药、运动、饮食）带来，人们会预期，一个不提及这些干预的反事实 $Y^{a}$ 没有良好定义，因而血压对卒中的因果效应定义不良。可是专家们一致同意“血压对卒中具有因果效应”，这是综合各类证据的结果（如把物理定律应用于血管、体外研究、动物实验中的血压钳制、尸检研究、降低血压的干预研究）。

解决这一表面矛盾的一种方式是：把“血压对这个人卒中具有因果效应”这类陈述解释为一种信念——相信存在某种干预 $X$（可能尚未被发现），它能把该个体的 $A$ 值改变为 $a$，但除了通过 $A$ 之外，对该个体的结局 $Y$ 没有（直接）效应。我们把这种干预称为对 $a$ 的\emph{有效无直接效应}（effective no-direct-effect，ENDE）干预，记作 ENDE($a$)。ENDE($a$) 干预的存在，在逻辑上等价于断言反事实 $Y^{a}$ 与联合反事实 $Y^{x,a}$ 都有良好定义且相等。

\mnote{文献中，“$X$ 除经由 $A$ 之外对 $Y$ 没有（直接）效应”与“$X$ 对 $Y$ 的效应完全由 $A$ 中介”这两种说法是同义的，可以互换。第 23 章讨论因果中介。}

相信 ENDE 干预存在，为“血压对这个人卒中具有因果效应”这类陈述提供了依据。对相关因果机制了解得相对清楚的非干预变量 $A$，这种信念可能更容易被接受。例如，在考虑血压对卒中的效应时，已知降压药、运动与饮食都能改变血压 $A$。其中，ENDE 干预可能是各种降压药，而不是饮食或运动——后两者已知会经由血压之外的途径影响卒中。在其他情形下，ENDE 干预的信念更难辩护，例如社会经济地位对心肌梗死的效应。

本节前面说过，对给定总体，平均因果效应 $\E[Y^{a}-Y^{a'}]=\Pr[Y^{a}=1]-\Pr[Y^{a'}=1]$ 只有在总体中每个个体的个体因果效应都有良好定义时才有良好定义。因此，对非干预 $A$，若总体中每个个体都存在 ENDE($a$) 与 ENDE($a'$) 干预，则总体中的平均因果效应有良好定义。如果对某些个体而言，没有任何 ENDE 干预能把 $A$ 的值改变为 $a$ 与 $a'$，那么总体中的因果效应就无从定义。也就是说，定义良好的总体效应要求研究者把总体限制在 ENDE 干预确实存在的那些个体上。见精要 3.4。

\mnote{为教学目的，我们忽略了血压 $A=(A_{\mathrm{sys}},A_{\mathrm{dia}})$ 具有收缩压与舒张压两个分量这一事实。若不忽略，那么“$X$ 除经由 $A$ 之外对 $Y$ 没有直接效应”这一陈述将断言反事实 $Y^{a_{\mathrm{sys}},a_{\mathrm{dia}}}$ 与联合反事实 $Y^{x,a_{\mathrm{sys}},a_{\mathrm{dia}}}$ 都有良好定义且相等。}

使 ENDE 干预的存在更可信的一种方式，是考虑 $A$ 值的微小变化。例如，考虑血压改变 $\Delta$ mm Hg 的效应，$\Delta$ 非常接近于 0。对（基本上）所有个体都更可能存在有效的干预；此外，直接效应可能只在考虑较大变化时才出现，例如在低剂量下没有直接效应的降压药，在高剂量下可能因毒性而产生直接效应。因此，对较小的 $\Delta$，所有个体的反事实 $Y^{\Delta}$ 以及平均因果效应 $\E[Y^{\Delta}-Y^{\Delta=0}]$ 都将（基本上）有良好定义。这里 $Y^{\Delta}$ 等于 $Y^{a}$ 在 $a=A-\Delta$ 处的取值，其中 $A$ 是个体在 ENDE 干预开始前的血压。

在专家们对非干预变量 $A$ 是否有相应的、定义良好的反事实 $Y^{a}$ 意见不一的情形下，也许可以设计一项研究，来支持“干预 $X$ 对 $Y$ 的效应完全由 $A$ 中介”这一命题。在精要 6.4 中，引入因果图之后，我们讨论这样一种研究设计。

\renewcommand{\FPname}{精要（Fine Point）}
\begin{finepoint}[干预主义路径的因果推断]
本书描述的因果推断框架，建立在定义良好的反事实结局之存在上，而反事实结局又由定义良好的干预来定义。我们的框架有时被称为“干预主义”框架，它提供了形式化的工具与语言，以便就干预的因果推断展开精确讨论，并解释数据分析所产生的数值量。

然而，关于因果推断的科学对话常常并不明确提及反事实结局或干预。科学文献中发表的许多论文涉及那些不包含可辨认干预的因果问题，却含有被解释为因果效应度量的数值量。因此，本书提出干预主义路径的一个扩展，通过澄清对非干预作有意义的定量因果推断所需的条件，把干预与非干预的因果推断统一起来。我们希望这一扩展能帮助许多学科的科学家，尤其是健康与社会科学领域。

我们并不主张干预主义路径是唯一有效的因果性哲学。然而据我们所知，确实不存在其他框架能给出具有实际可解释性的因果效应估计。
\end{finepoint}
\renewcommand{\FPname}{精要}

\section{一致性：其次，把反事实与观测数据联系起来}

提醒一下，一致性条件说的是：对所有 $A=a$ 的个体，$Y^{a}=Y^{A}=Y$。上一节讨论了一致性的第一个组成部分：足够定义良好的反事实结局 $Y^{a}$，不再残留有意义的模糊性。本节讨论一致性的第二个组成部分：把反事实结局与观测结局联系起来，即对 $A=a$ 的个体，$Y^{a}=Y$ 中的那个“等号”。

当干预 $a$ 确实由研究者按实验方案实施时，反事实与观测结局之间的联系毫无争议。对接受治疗取值 $A=a$ 的个体，观测结局 $Y=Y^{A}$ 按定义就等于反事实结局 $Y^{a}$。例如，在上一节的随机化实验中，一致性在每项实验所实施的那个版本的“心脏移植”之下成立。

当干预确实在真实世界中实施过（即便不是由研究者实施）时，类似的推理也适用于观察性研究。设想我们收集了符合移植条件的心脏病个体的数据，他们在某一时间、某一地点作为其医疗照护的一部分，被提供心脏移植（$A=1$）或药物治疗（$A=0$）。“心脏移植”这一干预 $a=1$ 的定义，对应于该时该地提供心脏移植之后所发生的一切程序；“药物治疗”$a=0$ 同理。对研究中 $A=1$ 的个体，在提供心脏移植之下的反事实结局 $Y^{a=1}$ 等于她的观测结局 $Y=Y^{A}$。当然，如果研究者无法大致刻画“心脏移植”$a=1$ 在该情境中的含义，这项观察性研究所针对的因果效应对其他总体的相关性就值得怀疑。

\mnote{关于使用观察性医疗数据库研究心脏移植效应时实际问题的扩展讨论，见 Macenati et al. (2024)。}

因此，用观察性数据时，研究中选择何种干预取决于可用的数据。例如，设想某观察性研究的研究者仔细定义了一个干预“心脏移植”$a=1$，精确规定了术前程序、麻醉、手术技术、术后程序与免疫抑制治疗。然而，数据中关于心脏移植的唯一信息，是一个人是否接受了心脏移植的指示变量 $B$。那么，对 $B=1$ 的个体，定义良好的反事实结局 $Y^{a=1}$ 未必等于该个体的观测结局 $Y$。

注意，我们对（假想的）干预 $a=1$ 与数据中（观测到的）变量 $B$ 使用了不同的字母。由于一致性条件说的是对所有 $A=a$ 的个体 $Y^{a}=Y^{A}=Y$，用同一个字母表示观测变量与假想干预，只保留给“每个个体的观测值 $A=a$ 恰好对应于定义良好的干预 $a$”的情形。换言之，能够描述一个定义良好的干预 $a$，并不足以实现一致性。我们还需要能够把定义良好的反事实结局 $Y^{a}$ 与观测结局 $Y$ 联系起来。

\renewcommand{\FPname}{精要（Fine Point）}
\begin{finepoint}[归因分数]
我们已经描述过一些效应度量，如因果风险比 $\Pr[Y^{a=1}=1]/\Pr[Y^{a=0}=1]$ 与因果风险差 $\Pr[Y^{a=1}=1]-\Pr[Y^{a=0}=1]$，它们把治疗 $a=1$ 之下的反事实风险与治疗 $a=0$ 之下的反事实风险相比较。然而，人们也可能关心把观测风险与治疗 $a=1$ 或 $a=0$ 之下的反事实风险相比较的度量。后一种对比使我们能在观察性研究中计算可归因于治疗的病例比例，即“若未发生治疗便不会发生的病例”所占的比例。例如，设想我们总体中的全部 20 人都参加了一场晚宴，席上供应的是神食（$A=1$）或琼浆（$A=0$）。第二天，接受 $A=1$ 的 10 人中有 7 人、接受 $A=0$ 的 10 人中有 1 人患病。为简单起见，假定受治疗者与未受治疗者可交换，从而因果风险比为 $0.7/0.1=7$，因果风险差为 $0.7-0.1=0.6$。（在条件随机化实验中，人们会通过标准化或逆概率加权计算这些效应度量。）后来发现，神食被一群鸽子污染了，这解释了因果风险比与因果风险差所概括的风险升高。现在我们要回答的问题是：“有多大比例的病例可归因于食用神食？”

在这项研究中，我们观察到 8 例病例，即观测风险为 $\Pr[Y=1]=8/20=0.4$。若人人都接受 $a=0$，本会观测到的风险为 $\Pr[Y^{a=0}=1]=0.1$。这两个风险之差为 $0.4-0.1=0.3$。也就是说，多出 30\% 的个体确实患病，但若总体中人人都接受 $a=0$ 而非其实际治疗 $A$，他们本不会患病。由于 $0.3/0.4=0.75$，我们说 75\% 的病例可归因于治疗 $a=1$：与观察到的 8 例相比，若人人都接受 $a=0$，本只会发生 2 例。这一超额分数或\emph{归因分数}定义为

\[
\frac{\Pr[Y=1]-\Pr[Y^{a=0}=1]}{\Pr[Y=1]}.
\]

关于超额分数在充分组分-原因框架中的讨论，见精要 5.4。

超额分数通常不同于\emph{病因分数}——归因分数的另一种版本，定义为由暴露在机制上引起的病例所占的比例。例如，设想未受治疗者（$A=0$）若接受治疗本会有 7 例病例，但这 7 例并不包含实际发生的那 1 例未受治疗病例，即治疗产生了 7 例却阻止了 1 例。再设想，受治疗者若未受治疗本只会发生 1 例，且该例不同于他们实际发生的 7 例。那么超额分数就不等于病因分数。此时超额分数是病因分数的一个下界。由于病因分数不依赖超额病例这一概念，它只能在强假设之下的随机化实验中计算。见 Greenland and Robins, 1988 与 Robins and Greenland, 1989。
\end{finepoint}
\renewcommand{\FPname}{精要}

反过来，某个变量 $A$ 恰好被记录在数据集中，并不保证 $A=a$ 可以联系到一个定义良好的干预 $a$。也就是说，测量了变量 $A$ 并不保证“$A$ 的因果效应”是一个有意义的概念，因为相应的干预定义不良，如前节所述。有定义良好的干预 $a$、却在数据中没有任何人 $A=a$，这一问题可视为非正值性的一种极端形式。

在观察性研究中实现一致性可能很困难。一个好做法是让我们的推理尽可能透明，以便他人能直接质疑我们关于一致性的论证以及我们对结果的解释。下一节描述一个提高这种透明度的程序。精要 3.5 说明如何在一致性之下，用观察性数据计算可归因于治疗的病例比例。

\section{目标试验}

本章探讨了三个条件——可交换性、正值性、一致性——它们有助于把观察性研究等同于条件随机化实验。因此，当研究者假定这三个条件成立时，他们的观察性分析就可以被视为在尝试模拟某个（假想的）随机化实验，该实验将量化所关心的平均因果效应。我们把那个假想实验称为\emph{目标实验}或\emph{目标试验}。

\mnote{目标试验——或其逻辑等价物——长期以来一直是因果推断框架的核心。Dorn (1953)、Wold (1954)、Cochran (1972)、Rubin (1974)、Feinstein (1971) 与 Dawid (2000) 都使用过这一概念。Robins (1986) 把它推广到随时间变化的处理。}

对每一个我们想用观察性数据估计的因果效应，我们可以 (i) 规定我们想要但无法实施的目标试验之方案，(ii) 描述如何用观察性数据来模拟该目标试验。如果模拟成功，那么观察性研究的结果与目标试验（若曾实施）的结果之间就没有差别。

试验方案的关键组成部分是：合格标准、干预（或一般地说，治疗策略）、分配、结局、随访的起点与终点，以及因果对比。合格标准的选择必须使所有个体原则上都能接受所有所关心的治疗（否则平均因果效应无从定义，因为某些个体因果效应无从定义）。为简单起见，在本书第一部分与第二部分中，我们只考虑每个个体仅在一个时点满足合格标准的情形（该时点成为随访起点，即时间零点）。关于目标试验框架的更多讨论见第 22 章。

\mnote{目标试验的关键组成部分由 Hernán and Robins (2016) 与 Hernán et al. (2025) 给出。缩写 PICO（Population、Intervention、Comparator、Outcome）有时用来概括其中若干组成部分（Richardson et al. 1995）。}

规定目标试验，是定义反事实结局 $Y^{a}$、从而定义所关心因果效应的自然方式；下面我们分别就实验与观察性研究加以讨论。

在由研究者控制治疗取值 $A$ 的实验中，反事实结局 $Y^{a}$ 有良好定义，因为治疗变量 $A$ 的每个可能取值 $a$ 都指向研究者施加于研究单元的一个定义良好的干预（行动）。例如，在一项关于低剂量阿司匹林 $A$ 与卒中 $Y$ 的随机化试验中，若参与者完全遵守所分配的治疗，则变量 $A$ 代表研究者按试验方案向参与者分配阿司匹林（$A=1$）或不分配阿司匹林（$A=0$）的行动。如 3.4 节所述，试验总体中每个人的反事实结局 $Y^{a}$ 都有良好定义，这意味着因果效应 $\Pr[Y^{a=1}=1]-\Pr[Y^{a=0}=1]$ 有良好定义。另一种表述是：该因果效应可以用一个目标试验来表达，而在这个例子里，目标试验恰好就是研究者实际实施的那项试验。

在观察性研究中，如果变量 $A$ 的每个可能取值 $a$ 都对应于某个主体（不必是研究者）本可施加于研究单元的定义良好的干预（行动），那么反事实结局 $Y^{a}$ 同样有良好定义。例如，考虑一个观察性数据库，其中 $A$ 是是否开出阿司匹林的指示变量。在限于“对阿司匹林既无绝对适应证也无绝对禁忌证”的个体的总体中，多数专家会把“被开出阿司匹林”与“未被开出阿司匹林”都视为足够定义良好的干预。也就是说，尽管并未实际实施随机化试验，专家们倾向于同意这样一项（目标）试验本可以进行。因此，对所关心因果效应的精确刻画，等价于规定该目标试验的方案。

现在考虑这样的观察性研究：$A$ 所能取的取值 $a$ 并不对应于假想目标试验中定义良好的干预。考虑三个这样的非干预变量 $A$：基因变异、血压（前几节已讨论过）与社会经济地位。目标试验框架能否用于这种情形？

\mnote{判断一个变量 $A$ 是否对应干预，常常有争议。有些变量毫无疑问是干预（随机化试验中的治疗分配），另一些则毫无疑问是非干预（社会经济地位）；其间存在渐变，专家们可能对某个变量是否为干预意见不一。}

第一，设想变量 $A$ 描述个体的基因型。按 3.4 节的讨论，若存在 ENDE 干预 $X$，可在受孕时把目标总体中每个个体的基因型改变为 $a$，而除经由 $A$ 值改变之外不影响 $Y$，则反事实结局 $Y^{a}$ 有良好定义。也就是说，对所关心的目标总体中每个个体，定义良好的反事实 $Y^{a}$ 与 $Y^{x,a}$ 相等。如果这些 ENDE 干预 $X$ 存在，那么因果效应有良好定义，并可通过规定一项目标试验的方案来精确刻画——该试验把个体随机分配到 $X$ 而不是 $A$。因此，任何关于“$A$ 的因果效应是否有良好定义”的争论，我们都应视为关于“专家们是否相信针对 ENDE 干预 $X$ 的目标试验可行”的争论。注意，干预 $X$ 不必是已知的。事实上，几十年前，多数专家认为基因型 $A$ 的因果效应有良好定义，因为他们相信某种尚待发现的技术 $X$ 可以改变 $A$ 的值而不直接影响远端结局。果然，基因编辑技术最终被开发出来了。

第二，设想变量 $A$ 表示个体的收缩压（单位 mm Hg）。同样，若存在 ENDE 干预 $X$，可把目标总体中每个个体的血压改变为 $a$，而除经由 $A$ 值改变之外不影响 $Y$，则对总体中所有个体反事实结局 $Y^{a}$ 都有良好定义。如 3.4 节所述，许多专家相信某些药物干预在某些目标总体中满足这些标准，尽管也许只在收缩压变化的有限范围内成立。因此，专家们可以通过规定一项目标试验的方案，来定义特定总体中 $A$ 的适度差异的因果效应——该试验把该总体中的个体随机分配到 ENDE 干预 $X$。

第三，设想变量 $A$ 表示社会经济地位的水平。多数专家怀疑存在这样的干预 $X$：它把个体置于某一社会经济地位 $a$，而不通过其他途径影响结局 $Y$（事实上，专家们甚至可能对“社会经济地位”指什么意见不一）。若不相信 ENDE 干预存在，社会经济地位的因果效应就不能被视为定义良好，等价地说，就无法规定目标试验。另见精要 3.6。

\mnote{其他作者认为，在许多专家不承认 ENDE 干预 $X$ 存在的情形下，“$A$ 对 $Y$ 的因果效应”仍有良好定义。见 Pearl (2009)、Schwartz et al. (2016)、Glymour and Spiegelman (2016)。}

总之，若 $A$ 表示一个干预，则目标试验是把满足合格标准的个体随机分配到不同 $a$ 值的试验；若 $A$ 表示非干预，且我们相信 ENDE 干预 $X$ 存在，则目标试验是把满足合格标准的个体随机分配到某个干预 $X$ 以把 $A$ 的值改变为 $a$ 的试验。换言之，对非干预 $A$ 的定量因果推断，要求假定 ENDE 干预 $X$ 存在，即便我们无法具体指明它们。于是 $A$ 的效应可以等同于 $X$ 的效应。没有这一假定，反事实结局 $Y^{a}$ 就没有良好定义，$A$ 的因果效应也没有良好定义。

一旦规定了目标试验方案，有效的模拟就需要足够的信息，以便把方案的各个要素映射到观察性数据集上，例如识别合格个体、按所关心的干预把他们分类、并在随访期间确定其结局。此外，在用上一章描述的方法——逆概率加权或标准化——计算因果效应时，有效的目标试验模拟还需要一个含有足够调整变量 $L$ 的数据集，以实现治疗组之间的条件可交换性。例如，对“血压 $A$ 对卒中 $Y$ 的效应”感兴趣并相信 ENDE 干预 $X$ 存在的研究者，需要调整那些能使不同血压水平的个体条件可交换的变量。由于许多研究者相信饮食与运动同时影响体重 $A$ 与卒中风险 $Y$，他们将需要调整这些变量，以期实现条件可交换性。在本书后面（见第 16 章），我们考虑模拟目标试验的另一些识别条件。

\renewcommand{\FPname}{精要（Fine Point）}
\begin{finepoint}[非干预的因果推断何时不受质疑]
如果专家们一致认为变量 $A$ 是一个干预，那就无须讨论 ENDE 干预 $X$ 是否存在。反之，当 $A$ 是非干预时，关于 $A$ 的定量因果推断就要求研究者讨论 ENDE 干预的可能性。然而，在某些情形下，我们照例谈论非干预 $A$ 的因果效应，却不明确讨论 ENDE 干预。这种情况出现在专家们对 ENDE 干预的存在并无疑问之时，如下例所示。

考虑一项随机化试验，个体被分配到治疗（$Z=1$）或无治疗（$Z=0$），同时我们收集他们实际接受的治疗 $A$（1 或 0）的数据。分配 $Z$ 的效应（即所谓意向性治疗效应）有良好定义，因为 $Z$ 是一个干预；但 $A$ 的效应又如何呢？$A$ 并不是干预。

如果试验中完全遵守治疗分配，则 $Z=A$，因为每个被分配到某一治疗取值的个体实际都接受了该治疗取值。那么，若 $Z$ 没有直接效应，$Z$ 就是 $A$ 的一个 ENDE 干预，我们无须区分 $Z$ 与 $A$。在本章正文的例子中，我们没有区分分配的治疗与实际接受的治疗，因为我们假定了完全遵守以及分配没有直接效应。

如果试验中并非完全遵守，则某些个体的 $Z$ 与 $A$ 取值不同。也就是说，$Z$ 不是对所有个体都构成 ENDE 干预，因为即便假定 $Z$ 对 $Y$ 没有直接效应，它也并不能“有效”地把每个个体的 $A$ 设为 1 或 0。无论如何，$Z$ 的效应仍有良好定义，因为它是干预；但 $A$（一个非干预）的效应现在只有在假定存在针对 $A$ 的 ENDE 干预 $X$ 时才有良好定义。在此语境下，多数专家会同意这类 ENDE 干预存在，例如在个体不知情的情况下把阿司匹林加入其食物。事实上，这些 ENDE 干预的存在似乎不言自明，以致我们在随机化试验中谈论“实际接受治疗 $A$ 的效应”而实际意指“$X$ 的效应”时，从不明确讨论它们。

尽管如此，无论 ENDE 干预 $X$ 是否存在，总体中总有这样一个子集：$Z$ 对他们有效——即那些若 $Z=1$ 则会有 $A=1$、若 $Z=0$ 则会有 $A=0$ 的个体。在这个子集中，即便 $Z$ 对总体中其他个体不构成 ENDE 干预，$Z$ 也是 ENDE 干预。因此，在这个子集中，$A$ 的效应有良好定义且等于 $Z$ 的效应。在第 16 章，我们把这个子集称为“依从者”，并说明尽管无法识别哪些个体属于该子集，但在附加假设之下可以识别依从者中 $A$ 的效应。

最后，如果 $Z$ 对 $Y$ 有不由 $A$ 传递的直接效应，那么 $Z$ 对 $Y$ 的效应将不同于 $X$ 对 $Y$ 的效应，即便在完全遵守的情况下也是如此。
\end{finepoint}
\renewcommand{\FPname}{精要}

把观察性分析锚定在目标试验上，能让因果推断对决策者——政策制定者、临床医生、监管者，还有你——更有相关性，帮助他们识别那些隐含极端或不可能干预的数据分析。更多讨论见下文与精要 3.7。

\renewcommand{\FPname}{精要（Fine Point）}
\begin{finepoint}[目标试验模拟的一些局限]
全书以心脏移植对死亡率的平均因果效应为例。然而，这一效应可以有不同解释。例如，“心脏移植”治疗之下的平均反事实，可以解释为总体中 $n$ 个合格个体在下列两种干预之下的反事实结局之平均：(i) 所有 $n$ 个个体同时接受治疗的干预；或 (ii) 每个个体 $i$ 接受治疗、而其余 $n-1$ 个个体接受他们实际所接受治疗的干预。解释 (ii) 是 $n$ 个干预的平均，对必须决定病人是否应接受心脏移植的医患双方高度相关。这也是本书一直隐含使用的解释。解释 (i) 所涉及的干预没有良好定义，因为它没有说明为了增加心脏供给、提高实施全部移植的能力，究竟会采用重设计医疗系统的诸多方式中的哪一种。

如果我们精确规定了医疗系统的重设计，那么我们现有的观察性数据就不够用，因为这些数据是在一个不包含那些改变的医疗系统下产生的。例如，某个医疗系统可能通过对移植器官质量降低筛选标准，来为所有合格个体提供心脏移植，而这会影响心脏移植之下的反事实结局。观察性数据可能不足以刻画把某项干预扩大规模、在全系统实施时的效应。

这一讨论与干扰有关（精要 1.1）。然而，与这里强调的问题不同，干扰文献一般假定解释 (i) 之下的反事实结局有良好定义（因为无须规定任何结构性的系统改变）。Hernán et al. (2025) 还描述了其他一些例子，其中目标试验的组成部分无法直接映射到观察性数据。
\end{finepoint}
\renewcommand{\FPname}{精要}

例如，设想我们有体重数据，并决定做一项分析，比较随访开始时不同体重水平（如体质指数 25 与体质指数 30）者的死亡风险。这一比较意味着模拟这样一项目标试验：个体在随访开始时被分配到一个瞬时的、且可能非常大的体重变化。由于我们不知道有任何 ENDE 干预能让人瞬时减掉这么多体重，我们可以提出对数据分析的修改，使其隐含一项经过修改的、干预更合理的目标试验，例如：某些个体被分配为从 40 岁起每年降低 5\% 的体质指数，只要其体质指数仍高于 25 就一直持续。（本书第三部分围绕的就是像这样随时间持续下去的干预。）

专家们也许仍对“能否实施这种经过修改的、体重逐步变化的目标试验中的 ENDE 干预”意见不一，但他们都会同意：这类干预比瞬时减重的干预更可能存在。因此，一种可行的做法是：假定这些更现实的干预存在，继续进行包含预后因素 $L$ 调整的分析，并接受其结果是不可作因果解释的、也可能不足以指导决策。好的一面是，我们的分析或许能提供一些关于 ENDE 干预（若存在）之效应的信息，或启发未来开发使这些干预得以实现的技术。

\mnote{在研究减重与心脏病时，Danaei et al. (2016) 考虑了会导致多年内体重逐步下降的未知干预。}

至少，规定目标试验有助于我们把因果问题对准那些可行、或至少并非完全不现实的干预。关心可操作因果推断的决策者，因而可能会把目标试验框架视为一个自然的起点。

\renewcommand{\TPname}{技术注记（Technical Point）}
\begin{techpoint}[递归替换]
给定一组按时间排序的变量，\emph{一步前瞻反事实}是指：在所有可被干预的先前变量都已被干预的情况下，某变量的反事实取值。设有按时间顺序排列的变量 $L,A,M,Y$，且对 $L,A,M$ 的干预都有良好定义。那么反事实 $L$、$A^{l}$、$M^{l,a}$、$Y^{l,a,m}$ 就是一步前瞻反事实。另一方面，若对 $L$ 的干预没有良好定义，则一步前瞻反事实变为 $L$、$A$、$M^{a}$、$Y^{a,m}$。其他一切事实量与定义良好的反事实，都可以通过“递归替换”由一步前瞻反事实构造出来（即是它们的函数）。有了一步前瞻反事实 $L$、$A^{l}$、$M^{l,a}$、$Y^{l,a,m}$，我们可以如下使用递归替换：

\begin{itemize}
\item $A=A^{L}$：把一步前瞻反事实 $A^{l}$ 在观测到的 $L$ 处取值。
\item $M^{a}=M^{L,a}$：把一步前瞻反事实 $M^{l,a}$ 在观测到的 $L$ 处取值。
\item $M=M^{L,A}=M^{L,A^{L}}$：把 $M^{l,a}$ 在观测到的 $L$ 与 $A$ 处取值。
\item $M^{l}=M^{l,A^{l}}$：把反事实 $M^{l,a}$ 在 $l$ 与反事实 $A^{l}$ 处取值。
\item $Y^{a}=Y^{L,a,M^{a}}=Y^{L,a,M^{L,a}}$：把反事实 $Y^{l,a,m}$ 在 $L$、$a$ 与 $M^{a}$ 处取值。
\item $Y^{m}=Y^{L,A,m}=Y^{L,A^{L},m}$。
\item $Y^{l}=Y^{l,A^{l},M^{l}}=Y^{l,A^{l},M^{l,A^{l}}}$。
\item $Y^{l,m}=Y^{l,A^{l},m}$。
\item $Y=Y^{L,A,M}=Y^{L,A^{L},M^{L,A^{L}}}$。
\end{itemize}

一步前瞻反事实也编码了无直接效应的假设。例如，若 $a$ 对 $Y$ 没有直接效应（相对于 $L$ 与 $M$），则 $Y^{l,a,m}=Y^{l,m}=Y^{l,A^{l},m}$，从而一步前瞻反事实变为 $L$、$A^{l}$、$M^{l,a}$、$Y^{l,m}$。

现在回到正文所述的心脏移植实验。在第一个实验中，\mbox{$Y^{a_0=1,a_1^{*},a_2^{*},a_3^{*}}$} 是在分配到心脏移植（$a_0=1$）、且干预各组成部分被详细规定为 $a_1^{*}$、$a_2^{*}$、$a_3^{*}$ 之下的反事实结局。令 $\mathcal{A}_1,\mathcal{A}_2,\mathcal{A}_3$ 为可能的、定义良好的 $(a_1,a_2,a_3)$ 干预之集合，其中包含 $(a_1^{*},a_2^{*},a_3^{*})$。一步前瞻反事实是 $A_1^{a_0=1}$、$A_2^{a_0=1,a_1}$、$A_3^{a_0=1,a_1,a_2}$ 与 $Y^{a_0=1,a_1,a_2,a_3}$。在第二个实验中，\mbox{$Y^{a_0=1}=Y^{a_0=1,A_1^{a_0=1},A_2^{a_0=1},A_3^{a_0=1}}$} 是在分配到心脏移植（$a_0=1$）之下的反事实结局，其后各组成部分取自然值
\[
A_1^{a_0=1}=A_1^{A_0=1},\qquad
A_2^{a_0=1}=A_2^{a_0=1,A_1^{a_0=1}},\qquad
A_3^{a_0=1}=A_3^{a_0=1,A_1^{a_0=1},A_2^{a_0=1}},
\]
全部用一步前瞻反事实写出。注意，当除分配到心脏移植 $a_0=1$ 之外未规定任何干预时，$A_1^{a_0=1}$、$A_2^{a_0=1}$ 与 $A_3^{a_0=1}$ 是随机的（即因人而异）反事实，取值于 $\mathcal{A}_1,\mathcal{A}_2,\mathcal{A}_3$。

递归替换说明了这一点。一般而言，把第二个试验的结局分布 \mbox{$Y^{a_0=1}=Y^{a_0=1,A_1^{a_0=1},A_2^{a_0=1},A_3^{a_0=1}}$} 迁移到另一总体，可能比迁移第一个试验的结局分布 \mbox{$Y^{a_0=1,a_1^{*},a_2^{*},a_3^{*}}$} 更困难：若 $A_1^{a_0=1}$、$A_2^{a_0=1}$ 或 $A_3^{a_0=1}$ 中一个或多个的分布在两个总体之间不同，则 $Y^{a_0=1,a_1,a_2,a_3}$ 的分布通常也会在两个总体之间不同。

假定一步前瞻反事实有良好定义，上述递归替换同样适用于观察性研究与随机化研究，因为一步前瞻反事实的定义只涉及反事实与事实量之间的逻辑关系，与研究类型无关。
\end{techpoint}
\renewcommand{\TPname}{技术注记}
